% titel: Blatt 1, A1: Vektorrechnung
% skript: Kap. Grundlagen Vektorrechnung

% Betrag, Skalarprodukt, Kreuzprodukt am konkreten Zahlenpaar.

Die Aufgaben loest man mit den drei Standard-Operationen der
Vektorrechnung: Betrag, Skalarprodukt und Kreuzprodukt.

\begin{werkzeugkasten}[Werkzeuge]
\textbf{Betrag} eines Vektors $\vec a=(a_x,a_y,a_z)$:
\begin{equation}
  |\vec a|=\sqrt{a_x^2+a_y^2+a_z^2}\;.
\end{equation}

\textbf{Skalarprodukt} (Skalar, ergibt den Winkel):
\begin{equation}
  \vec a\cdot\vec b = a_xb_x+a_ya_y+a_zb_z
  = |\vec a|\,|\vec b|\cos\theta
  \;\Rightarrow\;
  \cos\theta=\frac{\vec a\cdot\vec b}{|\vec a|\,|\vec b|}\;.
\end{equation}

\textbf{Kreuzprodukt} (Vektor senkrecht auf beiden Faktoren):
\begin{equation}
  \vec c=\vec a\times\vec b
  =\begin{pmatrix}
     a_yb_z-a_zb_y\\ a_zb_x-a_xb_z\\ a_xb_y-a_yb_x
   \end{pmatrix},
  \qquad |\vec c|=|\vec a|\,|\vec b|\sin\theta\;.
\end{equation}

\textbf{Addition komponentenweise:} $\vec a+\vec b=(a_x+b_x,\dots)$.
\end{werkzeugkasten}

\begin{werkzeugkasten}[Vorgehen]
\begin{itemize}
  \item Einen Winkel erhaelt man immer ueber $\cos\theta$ aus dem
        Skalarprodukt -- nie ueber die Einzelne koordinaten raten.
  \item $|\vec c|$ ist zugleich der Flaeche des von $\vec a,\vec b$
        aufgespannten Parallelogramms proportional.
  \item Senkrechtheit testet man mit dem Skalarprodukt: zwei Vektoren
        stehen senkrecht genau dann, wenn ihr Skalarprodukt null ist.
\end{itemize}
\end{werkzeugkasten}

\begin{achtung}[Typische Fallen]
\begin{itemize}
  \item Skalarprodukt und Kreuzprodukt verwechseln: das eine liefert eine
        Zahl, das andere einen Vektor.
  \item Beim Kreuzprodukt die Reihenfolge und die Vorzeichen der
        Kofaktoren (Muster $+,-,+$ in der ersten Zeile) beachten.
  \item Fuer den Winkel den Definitionsbereich von $\arccos$ bedenken;
        er liegt zwischen $0$ und $\pi$.
\end{itemize}
\end{achtung}
